% Additive incorporation. Does not modify the sources or the 211-page PDF.
% Provenance, dependencies, and scope in preservation_map.json.
\section{Grassmannian representation of readers and forms valued in the action line}
\label{sec:compat-lectores-formas}

Positive realizations preserve information that can subsequently be
used by physical readers. This preservation has a precise
formulation: a marked representation of the record allows its
coordinates to be reconstructed and every function using them to be transported by composition.
The geometry represents an outcome of the construction; its choice of
coordinates preserves the provenance of that outcome.

The starting domain consists of the states produced by APP,
TRIT, and TPK in the preceding sections. APP preserves its additive and
multiplicative sheets through remainder and quotient; TRIT selects the local
regime with its orientation; the composition
\(U_t=\operatorname{Upd}_t\circ\operatorname{Tra}_t\circ
\operatorname{Sel}_t\) updates position, record, and memory.
The prolongation
\[
 w_6\longrightarrow w_{12}\longrightarrow w_{18}\longrightarrow
 w_{24}\longrightarrow w_{30}\longrightarrow R_{36}\longrightarrow G_9
\]
preserves the prefix and distinguishes phase return from state return.
The spectral, solenoidal, gauge, cohomological, and determinantal realizations
remain correlative operations of the joint discrete
structure of the continuum. Record extraction and regional
representations take place in that same domain. In what follows, all of them
are previously constructed HMT outputs.

\subsection{The marked collection and its inverse}

Let \(L\geq2\), and let \(K=(K_1,\ldots,K_L)\) be a positive record
with charge \(Q=\sum_jK_j\). In the dodecaphasic application,
\(L=12\) and \(Q=6263\). Define
\[
 d_j=K_j/Q,\qquad t_0=0,\qquad
 t_j=\sum_{a=1}^{j}d_a,\qquad
 B(K)=\begin{pmatrix}1&\cdots&1\\t_0&\cdots&t_L\end{pmatrix}.
\]
The collection \(\mathcal G_L^{\rm aff}\) is the image of these row
planes in \(\operatorname{Gr}_{>0}(2,L+1)\), retaining the
ordered columns \(0,\ldots,L\), the endpoint markings, and the charge
\(Q\). Affine normalization fixes \(t_0=0,t_L=1\).
We write \(\Delta_{ij}\) for the minor of columns \(i,j\)
of any representative of the plane.

\begin{proposition}[Inverse of the marked realization]
\label{prop:compat-inversa-marcada}
The map \(K\mapsto([B(K)],Q)\) is injective. Its inverse
on the image is given by
\begin{equation}
 K_j=Q\frac{\Delta_{j-1,j}}{\Delta_{0L}},
 \qquad
 t_j=\frac{\Delta_{0j}}{\Delta_{0L}},
 \quad 1\leq j\leq L.
 \label{eq:compat-inversa-plucker}
\end{equation}
The ratios are invariant under any invertible change of basis
of the two rows.
\end{proposition}
\begin{proof}
In the declared representative,
\(\Delta_{ij}=t_j-t_i\), and \(\Delta_{0L}=1\).
The formulas recover the consecutive differences and their sums.
If \(B\) is replaced by \(AB\), with \(A\) invertible,
all minors are multiplied by \(\det A\); their ratios
remain unchanged. The charge recovers the scale preceding
normalization. Therefore two records with equal marked data
have the same coordinates.
\end{proof}

The condition of membership in \(\mathcal G_L^{\rm aff}\) matters.
The consecutive ratios of an arbitrary positive plane may
fail the identity \(\sum_j\Delta_{j-1,j}=\Delta_{0L}\).
In the constructed collection this identity is telescoping. The inverse acts
on the marked image just defined.

\subsection{Reader algebra and changes of chart}

Let \(\mathcal D\) be the image of an enriched-state representation
of the form \(x\mapsto(K(x),\xi(x))\). The symbol \(\xi\)
collects exactly the remaining coordinates used by the reader:
regional representations, origin and orientation, route, signature, boundary
conditions, and dimensional sections, according to the object considered.
These coordinates are not assumed to be independent. Define
\[
 \mathcal F:\mathcal D\longrightarrow\widetilde{\mathcal G},
 \qquad (K,\xi)\longmapsto([B(K)],Q,\xi),
 \qquad \widetilde{\mathcal G}:=\mathcal F(\mathcal D).
\]
The preceding proposition provides the inverse \(\mathcal I\) of
\(\mathcal F\) on its image.

\begin{theorem}[Faithful representation of the algebra of retained readers]
\label{thm:compat-algebra-lectores}
For each reader \(\mathcal R:\mathcal D\to V\), with declared
codomain \(V\), there exists a unique geometric reader
\(\mathcal R^{\rm geom}:\widetilde{\mathcal G}\to V\) such that
\[
 \mathcal R^{\rm geom}=\mathcal R\circ\mathcal I,
 \qquad \mathcal R^{\rm geom}\circ\mathcal F=\mathcal R.
\]
For scalar readers, this correspondence preserves sums, products,
the unit, and conjugation. It is an isomorphism from the algebra of functions on
\(\mathcal D\) to the corresponding algebra on
\(\widetilde{\mathcal G}\).
\end{theorem}
\begin{proof}
The composition is defined because \(\mathcal I\mathcal F\) is
the identity. If two geometric readers have the same composition
with \(\mathcal F\), they agree on its image, which is all of
\(\widetilde{\mathcal G}\). The algebraic operations hold
pointwise; for example,
\((\mathcal R_1\mathcal R_2)\circ\mathcal I=
(\mathcal R_1\circ\mathcal I)(\mathcal R_2\circ\mathcal I)\).
The inverse composition is restriction along \(\mathcal F\).
\end{proof}

Faithfulness refers to the representation \((K,\xi)\). For an
observable \(\mathcal O\) of the full enriched state, the descent
criterion is that \(\mathcal O(x)=\mathcal O(x')\) whenever
\((K(x),\xi(x))=(K(x'),\xi(x'))\). Necessity follows by
composition; sufficiency allows the value on a fibre to be defined
using any of its representatives. This criterion explicitly
preserves which information must accompany the record.

Now let \(\mathcal T\) be a triangulation of the labelled polygon
determining a positive Pl\"ucker chart. Its boundary coordinates
and a homogeneous normalization are included, for example
\(\Delta_{0L}=1\). If the diagonal \(ac\) is exchanged for
\(bd\), with \(a<b<c<d\), the change of chart is
\begin{equation}
 \Delta_{bd}=
 \frac{\Delta_{ab}\Delta_{cd}+\Delta_{ad}\Delta_{bc}}{\Delta_{ac}}.
 \label{eq:compat-cambio-carta}
\end{equation}
The denominator is positive, and the Pl\"ucker relation proves that both
charts evaluate the same point.

\begin{corollary}[Compatibility of readers with mutations]
\label{cor:compat-mutaciones}
Let \(r_{\mathcal T}\) and \(r_{\mathcal T'}\) be the expressions
of \(\mathcal R^{\rm geom}\) in two charts preserving the markings,
and let \(\mu_{\mathcal T\mathcal T'}\) be their positive transition. Then
\[
 r_{\mathcal T'}\circ\mu_{\mathcal T\mathcal T'}
 =r_{\mathcal T}
\]
on the common domain. A closed sequence of chart changes that
preserves the marked point preserves all these readers.
\end{corollary}
\begin{proof}
The two reconstructions through
\eqref{eq:compat-inversa-plucker} recover the same \(K\).
The coordinates \(\xi\) are transported with their markings. Applying
\(\mathcal R\) proves the equality; iterating it proves the claim
about a closed sequence.
\end{proof}

Readers can be expressed in different charts while remaining
functions of the same object. A permutation of labels acts
additionally on the positional marking; a transformation changing
the point acts on the reader's argument. Each operation has
its own transport law.

The diagonal action on columns provides an explicit check.
For \(\lambda_i>0\),
\(\Delta_{ij}\mapsto\lambda_i\lambda_j\Delta_{ij}\).
The cross-ratio \(\Delta_{ab}\Delta_{cd}/
(\Delta_{ad}\Delta_{bc})\) is preserved, while
\(\Delta_{j-1,j}/\Delta_{0L}\) acquires the factor
\(\lambda_{j-1}\lambda_j/(\lambda_0\lambda_L)\).
With \(t=(0,1/4,1/2,1)\) and \(\lambda=(1,2,1,1)\), the
reconstructed gaps would be \((1/2,1/2,1/2)\), whose sum
is \(3/2\). Cross-ratios alone preserve a quotient different
from the one required by the marked inverse.

\subsection{Explicit composition with action and constitutive response}

In the dodecaphasic chart, we write \(b=1000\) and
\begin{equation}
 \kappa^{\rm geom}=
 \frac{Q}{b^{12}-1}\sum_{j=1}^{12}
 b^{12-j}\frac{\Delta_{j-1,j}}{\Delta_{0,12}},
 \qquad
 \alpha^{\rm geom}=\pi_{\rm HMT}+e_{\rm HMT}
 -\varphi_{\rm HMT}-4-\kappa^{\rm geom}.
 \label{eq:compat-alpha-geometrica}
\end{equation}
These are the geometric expressions of the closure reader
\eqref{eq:mass-k-alpha}. The three regional representations and \(K\)
share a coinductive generation; the expression uses their internal
constructive order. The action reader
\eqref{eq:mass-k-action}, evaluated at \(\alpha^{\rm geom}\),
determines \(\eta_{\rm ret}^{\rm geom}\), and therefore
\[
 \hbar_{\rm ret}^{\rm geom}
 =\eta_{\rm ret}^{\rm geom}10^{-34}\mathcal U_S
 =\hbar_{\rm ret}.
\]
The equality holds in the action line \(\mathcal L_S\)
with its section \(\mathcal U_S\); it preserves the decade arising
from the determinantal reader \(\varphi_{\rm HMT}\alpha^{16}\).
The composition neither changes the section nor applies it a second time.

In the lifted angular chart preserving the origin and real branch
of the representation, the subsequent angular pair determines
\[
 x=\frac{\pi_{\rm HMT}}{180}(1000\alpha),\qquad
 y=\frac{\pi_{\rm HMT}}{90}(\eta_{\rm ret}+\alpha),\qquad
 q_\pm=e^{-x\mp y}.
\]
The branch is part of the retained markings: reduction of an angle
modulo \(360\) alone preserves different information. The equalities
below use the declared lift, in which the first
coordinate is \(1000\alpha\).
In the chamber \(x>|y|\), with fixed sheet projectors
\(P_\pm\), let \(T=q_+P_++q_-P_-\). The response of the corpus
compares incidences \(90\) and \(120\) through
\[
 \mathcal V=(I-T^{90})(I-T^{120})^{-1}
 =r_+P_++r_-P_-,\qquad
 r_\pm=f(q_\pm^{30}),\quad
 f(s)=\frac{1+s+s^2}{1+s+s^2+s^3}.
\]
The rule satisfies
\(\mathcal V\Sigma_4(T^{30})=\Sigma_3(T^{30})\), where
\(\Sigma_j(S)=\sum_{a=0}^{j-1}S^a\). Its uniqueness follows
from invertibility of \(\Sigma_4(T^{30})\).
The four normalized readings are
\[
 \widehat\varepsilon=r_+^2,\quad
 \widehat\mu=r_-^2,\quad
 \widehat Z=r_-/r_+,\quad
 \widehat c=(r_+r_-)^{-1}.
\]
In the dimensional lines of action, charge, and speed they are realized as
\[
 Z=\widehat Z\,u_Su_Q^{-2},\quad c=\widehat c\,u_C,
 \qquad
 \mu=\widehat\mu\,u_Su_C^{-1}u_Q^{-2},\quad
 \varepsilon=\widehat\varepsilon\,u_S^{-1}u_C^{-1}u_Q^2.
\]
The identities \(\mu\varepsilon=c^{-2}\) and
\(\mu/\varepsilon=Z^2\) follow by multiplication and division.
The chart-change corollary preserves this complete response,
with the sheet markings and dimensional sections included in
\(\xi\). The same argument applies to a mass reader with its
route, signature, and return operator retained, on the domain of
Section~\ref{sec:k-mass-bridge}.

\begin{theorem}[Recovery of the record from the labelled response]
\label{thm:compat-recuperacion-constitutiva}
With the regional representations, sheet orientation,
base \(b=1000\), origin, length twelve, and declared angular lift
fixed, the map
\(K\mapsto(\widehat\varepsilon,\widehat\mu)\) is injective
on the integer records of its physical image satisfying
\(0\leq K_j<b\) and the stated contractive chamber. Its exact inverse
is the composition of the following operations:
\begin{align*}
 r_+&=\sqrt{\widehat\varepsilon},&
 r_-&=\sqrt{\widehat\mu},\\
 r_\pm s_\pm^3+(r_\pm-1)(1+s_\pm+s_\pm^2)&=0,
 &0<s_\pm<1,\\
 x&=-\frac1{60}\log(s_+s_-),&
 \alpha&=\frac{180x}{1000\pi_{\rm HMT}},\\
 \kappa&=\pi_{\rm HMT}+e_{\rm HMT}-\varphi_{\rm HMT}-4-\alpha,
 &N&=(b^{12}-1)\kappa,\\
 K_j&=\left\lfloor N/b^{12-j}\right\rfloor\bmod b.
\end{align*}
\end{theorem}
\begin{proof}
The positive roots recover the labelled coefficients.
The derivative
\[
 f'(s)=-\frac{s^2(3+2s+s^2)}{(1+s+s^2+s^3)^2}<0
 \qquad(0<s<1)
\]
and its limits \(1\) and \(3/4\) prove that the cubic equation
has exactly one root in that interval. Since
\(s_\pm=e^{-30x\mp30y}\), their product is \(e^{-60x}\).
This recovers \(x\), \(\alpha\), and \(\kappa\).
On the image considered, \(N=\sum_jK_jb^{12-j}\) is an integer
between zero and \(b^{12}-1\). Successive Euclidean division uniquely recovers
its twelve digits, including leading zeros.
\end{proof}

The precision needed for numerical recovery can also be
expressed. With exact regional representations, an error
\(|\widetilde\alpha-\alpha|<1/[2(b^{12}-1)]\) implies
\(|\widetilde N-N|<1/2\); rounding to the nearest integer
recovers \(N\). If the regional representations have errors, the
same condition applies to the sum of errors in
\(\pi+e-\varphi-\alpha\). Exact injectivity and metrological
resolution are thus distinct properties.

The integer domain is essential. In an analytic prolongation of
positive records, perturbing three consecutive positions
proportionally to \((1,-(b+1),b)\) preserves both the sum and the
positional reading, since
\(1-(b+1)+b=0\) and \(b^2-b(b+1)+b=0\).
For sufficiently small perturbations it preserves positivity
and changes the Grassmannian point. The preceding injectivity uses
the condition of twelve integer digits; the continuous analytic collection
has different fibres.

\subsection{Refinement and inherited readers}

If each gap \(d_j\) is subdivided into positive gaps
\(d_{j,a}\) with \(\sum_ad_{j,a}=d_j\), the original endpoints
retain their positions. Let \(\rho\) be the operation of summing the successors
of each interval and \(H\) the inclusion of its endpoints in the new
ordered set. Then
\[
 \mathcal I_{\rm heredada}(B(\widetilde d),Q)
 =\bigl(Q\sum_ad_{j,a}\bigr)_j=K,
 \qquad
 \mathcal R_{\rm heredada}=\mathcal R(K,\xi).
\]
The equality proves commutativity between inherited reconstruction and
subdivision. In a sequence of subdivisions, sums associate;
the result persists in every finite composition. This preserves the
length-twelve reader and its markings. A reader defined on the
new blocks will have its own domain and comparison law.

\subsection{Canonical forms valued in a dimensional line}

Let \(P\) be one of the rank-one polytopes constructed in the
manuscript and \(P=\bigcup_\sigma P_\sigma\) an oriented
subdivision of the same support, with multiplicity one and compatible faces.
We write \(\Omega_\sigma\) and \(\Omega_P\) for their canonical
forms. Their identities are
\(\Omega_P=\sum_\sigma\Omega_\sigma\) and
\(\operatorname{Res}_F\Omega_P=\Omega_F\), with the orientation
convention fixed in the amplituhedral development.
The action line \(\mathcal L_S\) is taken to be constant over
this polytope; its section comes from the fixed HMT state.

\begin{theorem}[Transport of the form by an action section]
\label{thm:compat-forma-accion}
For \(a\in\mathcal L_S\), constant with respect to the coordinates
of the polytope,
\[
 \boldsymbol\Omega_P:=a\otimes\Omega_P
 =\sum_\sigma a\otimes\Omega_\sigma,
 \qquad
 \operatorname{Res}_F\boldsymbol\Omega_P=a\otimes\Omega_F.
\]
The identities are preserved under successive subdivisions and iterated
residues on flags of faces. They apply, in particular, to
\(a=\hbar_{\rm ret}\).
\end{theorem}
\begin{proof}
The tensor product with a line is linear. The residue acts on
the differential factor and commutes with the constant coefficient.
Compatibility of subdivisions and iterated residues is transported
by the same linear operations. Under a dimensional change of basis
\(u'_S=\lambda\,u_S\), with \(\lambda>0\), the numerical coefficient
changes by \(\lambda^{-1}\);
the tensor and its residues remain unchanged.
\end{proof}

\begin{theorem}[Gluing condition for boundary coefficients]
\label{thm:compat-pegado}
Let \(a_\sigma\) be sections of the same line, regular in a
neighbourhood of the faces and meromorphic in a common ambient realization.
Suppose their products with the canonical forms have poles of at
most first order along the faces considered. In the
relative interior of a common codimension-one face \(F\) of two cells
\(\sigma,\tau\), with opposite orientations,
\begin{equation}
 \operatorname{Res}_F\left(a_\sigma\Omega_\sigma+a_\tau\Omega_\tau\right)
 =\left(a_\sigma|_F-a_\tau|_F\right)\Omega_F.
 \label{eq:compat-defecto-frontera}
\end{equation}
Thus the pole of this pair on the divisor of \(F\) disappears
exactly when the two traces agree. If the coefficients are
constant and the cell adjacency graph is connected, cancellation
for each pair of cells across their common face is equivalent to the existence
of a single common coefficient \(a\).
\end{theorem}
\begin{proof}
With a transverse local coordinate \(z\), a logarithmic form is
written \(dz/z\wedge\omega+\eta\), where \(\eta\) is regular
with respect to that divisor. Regularity of \(a_\sigma\) gives
\(\operatorname{Res}_F(a_\sigma\Omega_\sigma)
=a_\sigma|_F\operatorname{Res}_F\Omega_\sigma\).
The two orientations produce \(\Omega_F\) and \(-\Omega_F\).
The form \(\Omega_F\) is nonzero in the relative interior of the
face; equality of residues is equivalent to equality of traces.
For a simple pole, vanishing of the residue removes the pole.
In the constant case, equality propagates along every
edge of the connected adjacency graph.
\end{proof}

For the global sum, all terms with a pole on the
same ambient divisor must be counted. If \(H\) is the hyperplane containing \(F\),
and \(I_H\) collects those terms, the general identity is
\[
 \operatorname{Res}_{H}\left(\sum_\nu a_\nu\Omega_\nu\right)
 =\sum_{\nu\in I_H}a_\nu|_H\operatorname{Res}_H\Omega_\nu,
\]
provided the coefficients are regular on \(H\). A coplanar
facet, even if it is not adjacent to the real face \(F\), contributes to this
sum. The two-term formula represents the global residue when
the remaining terms are regular on that divisor. Pairwise cancellation
of all interior facets removes their contributions;
the exterior facets preserve the corresponding boundary residues.

On a face of higher codimension, this formula is applied successively
with the induced orientations. The signs are those of cellular
incidence and satisfy cancellation of the boundary of a boundary. The collection of
local readers defines pairwise compatible gluing when their
traces agree after transport to the same coefficient line.
The global residue is obtained by summing the contributions of each divisor.
This statement gives an effective criterion for distributing a physical
reading over charts or cells without introducing spurious discontinuities.

\subsection{Pole cancellation and global normalization}

Interior cancellation determines a local condition. Identification
with a global canonical form also requires the pole divisor and
the normalized exterior residues. In projective space, two
rational top-degree forms with the same logarithmic poles
and the same residues differ by a holomorphic top-degree form;
\(H^0(\mathbb P^m,\Omega^m)=0\) proves their equality.
This is the global condition used in the polytope proof
\cite{positive}.

An interval example shows why variable coefficients
require this condition. For \(0<c<1\), let
\[
 \Omega_{0c}=\frac{c\,dx}{x(c-x)},\qquad
 \Omega_{c1}=\frac{(1-c)\,dx}{(x-c)(1-x)},\qquad
 \Omega_{01}=\frac{dx}{x(1-x)}.
\]
With \(a_1(x)=1+x(c-x)\) and \(a_2(x)=1\), their traces agree
at all relevant endpoints, and yet
\[
 a_1\Omega_{0c}+a_2\Omega_{c1}=\Omega_{01}+c\,dx.
\]
The additional term is regular in the affine chart and has a pole at
projective infinity. The residues at the finite endpoints and
cancellation at \(c\) are preserved, while the resulting form
has different global behaviour. Multiplying the example by a
section of \(\mathcal L_S\) gives the same dimensional check.

The overall consequence is precise: the marked Grassmannian
representation transports the retained readers; mutations preserve their
values when they change only the chart; refinement preserves their
inherited evaluations; and forms valued in the action line
are assembled through equality of traces and control of poles. Each
equality expresses the operation preserving the corresponding object.
